\chapter{Fraud Rings theory}
\label{ch:fraud_ring_theory}

\section{The Limitations of the Lonely Scammer}
Individual scammers typically rely on simple and intuitive techniques to introduce illicit funds into the legitimate economy. One of the most prevalent methods is known as ``smurfing''. In simple terms, smurfing involves breaking down a large sum of illicit money into many smaller, inconspicuous chunks to deposit them without raising suspicion. 

More formally, individual scammers rely on smurfing to avoid Anti-Money Laundering (AML) detection mechanisms, which usually flag transactions above a specific threshold $\tilde{x}$. By splitting their initial illicit funds $x$ into smaller transactions $x_n$, such that each transaction satisfies $x_n \leq \tilde{x}$, they minimize the risk of detection \cite{contreras2026strategic}. The mathematical expectation of profit $\pi_s$ for a scammer engaging in smurfing can be defined as:
\begin{equation}
    \pi_s = x(1-\nu)
\end{equation}
where $x$ represents the total initial illicit funds, and $\nu \in (0,1)$ is the friction cost or commission proportion associated with the smurfing process (e.g., the fee paid to mules). 

While smurfing is effective for small amounts, its scalability is severely limited. When attempting to launder much larger volumes of funds, an individual scammer operating from a single account faces a significantly higher probability of detection $(1-p)$, where $p$ is the probability of successfully avoiding detection. If caught, the scammer is subjected to a severe monetary penalty $T$, which includes both the total confiscation of the funds and associated judicial fines \cite{contreras2026strategic}. Therefore, operating as a ``lonely scammer'' becomes mathematically and practically unviable for large-scale operations due to the massive expected cost of $(1-p)T$.

\section{The Emergence of Fraud Rings}
To circumvent the limitations of individual operations and avoid getting caught, scammers evolved to form collusive networks, commonly referred to as ``fraud rings''. These rings operate by dispersing the illicit funds across a vast, interconnected network of accounts. They utilize a strategy of transaction layering, often characterized by a topological depth $L$, and network diversification to obscure the money trail \cite{contreras2026strategic, li2025scamsweeper}.

By spreading the funds through a complex topology, the fraud ring successfully increases its non-detection probability $p_L$, which now depends explicitly on the layering depth $L$. The expected profit of a fraud ring utilizing $L$ layers can be mathematically expressed as:
\begin{equation}
    E[\pi_L] = p_L x - (1-p_L)T - c(L)
\end{equation}
where $E[\pi_L]$ is the expected profit, $p_L$ is the enhanced probability of evading detection when $L$ layers are used, $x$ is the total illicit funds, $T$ is the total penalty if caught, and $c(L)$ represents the monotonic cost function associated with establishing and maintaining the $L$ transaction layers. Because the fraud ring dramatically increases the probability of evading detection ($p_L \to 1$ as $L$ increases), the expected payoff remains highly profitable even when factoring in the complex layering costs $c(L)$ \cite{contreras2026strategic}. 

\section{Graph Theory, GNNs, and Quantum Computing}
Because fraud rings create such complex, multi-layered transaction topologies, fraudulent behaviors can propagate with remarkable speed and sophistication \cite{awoyemi2017}. Traditional heuristic or rule-based detection tools---which often rely on simple random sampling or threshold rules---fail to identify them effectively. To counter these sophisticated rings, modern detection systems leverage Graph Theory and deep learning \cite{li2025scamsweeper}.

The transaction network is modeled as a directed graph $G = (V, E)$, where $V$ represents the set of accounts (nodes) and $E$ represents the set of transactions (edges). While classical Graph Neural Networks (GNNs) have proven highly successful in analyzing local neighborhood dependencies through message-passing mechanisms, they often fail to capture the higher-order, long-range dependencies critical for identifying distributed and coordinated fraudulent behaviors \cite{motie2024}. One technique to sample and study these dense, evolving networks is through temporal random walk subgraphs, which allow models to capture the dynamic sequence of transactions over time (this sampling mechanism will be detailed thoroughly in a subsequent implementation section).

In this work, we will utilize a Quantum Topological Graph Neural Network (QTGNN) to learn the dense, anomalous connectivity inherent to fraud rings. Drawing from recent advancements \cite{mainarticle}, QTGNN bridges quantum machine learning, graph theory, and topological data analysis, working exceptionally well for this task by addressing three fundamental challenges:

\begin{itemize}
    \item \textbf{Leveraging quantum embeddings for richer relational representation:} Traditional representations struggle to model the high-dimensional pathways of fraud rings. Quantum embeddings encode this complex relational information by exploiting quantum superposition and entanglement. This allows the model to represent multiple transaction pathways and interdependencies simultaneously, capturing subtle patterns distributed across the network that classical local message-passing GNNs often overlook \cite{mainarticle}.
    
    \item \textbf{Incorporating topological invariants for robustness and interpretability:} Topological Data Analysis (TDA) has gained significant attention for its ability to extract invariant features that characterize the global structure of data \cite{monkam2024}. By extracting topological descriptors such as persistent homology, the QTGNN captures global structural features that remain invariant to noise and minor perturbations. This not only makes the representations robust against the spurious connections generated by scammers but also enhances the model's interpretability by identifying the higher-order motifs driving the anomaly detection.
    
    \item \textbf{Ensuring scalability and practicality on NISQ devices:} While quantum models offer theoretical advantages, their practical deployment is heavily constrained by the limitations of Noisy Intermediate-Scale Quantum (NISQ) hardware \cite{umeda2019}. The QTGNN framework addresses these bottlenecks through a NISQ-aware design that employs circuit optimization, error mitigation, and hybrid quantum-classical computation strategies. This ensures that the quantum-enhanced graph learning remains feasible, scalable, and fully deployable on real-world financial transaction networks \cite{mainarticle}.
\end{itemize}

A comprehensive mathematical explanation of the QTGNN architecture and its implementation will be provided in a future section.